% =============================================================
%  Chapter — HVDC Topologies
%  Assets (SVG source + PDF build) live in this same subfolder.
% =============================================================
%\sectionslide{System Overview}{}


\begin{frame}{System configuration}
  \begin{columns}[c,onlytextwidth]
    \begin{column}{1\textwidth}
      \centering
        % \topoimg[width=\linewidth,height=0.85\textheight,keepaspectratio]{chapters/system-configuration/system_overview_poi_1}
        \topoimg[width=\linewidth,height=0.85\textheight,keepaspectratio]{chapters/system-configuration/system_overview_poi_2}
    \end{column}
  \end{columns}
\end{frame}


\begin{frame}{System configuration}
  \begin{itemize}
    \setlength{\itemsep}{0.5em}
    \item \textbf{Power} - Projects with $1000MW$ and $3000MW$ power transmission requirements
    \item \textbf{DC line} - 130 -200 miles, with priority to OHL, but cable sections can be implemented to accelerate the permission process.
    \item \textbf{Generation} - Geothermal power plants is the main generation source. Additional generation sources might be implemented once the transmission system is established.
    % \item \textbf{Load} - The main load are data centers. The interconnection voltage to the data centers to be defined ($36kV$ or $800V$).
    \item \textbf{Load} - The main load are data centers.
    \item \textbf{BESS} - Battery storage is assumed for backup and for datacenter load flicker compensation.
    \item \textbf{AC Chopper} - For fault ride through, installation of the AC chopper might be required.
  \end{itemize}
\end{frame}


\begin{frame}{System configuration}
  \begin{itemize}
    \setlength{\itemsep}{0.5em}
    \item Islanded operation with no interconnection to the AC grid
    \item Generation is build out in $300MW$ steps and is defined by the water availability \\ $\Rightarrow$ Staged buildout of the transmission line is beneficial
    \item High realiability requirements by datacenters
    % \item Execution time is critical
    \item Standardized hardware design approach to reduce hardware lead times and configuration variations
    \item Building up C\&P capabilities to reduce risk and execution times in project execution
    \item \textbf{Reduction of execution time is very critical}
    \item \textbf{4} initial projects, 2 x $1000MW$ and 2 x $3000GW$
  \end{itemize}
\end{frame}


\begin{frame}{HVDC Configuration - Dual SMP}
  \begin{columns}[c,onlytextwidth]
    \begin{column}{0.46\textwidth}
      \centering
      \topoimg[width=\linewidth]{chapters/system-configuration/dsmp}
    \end{column}
    \begin{column}{0.50\textwidth}
      \vspace{0.5em}
      \begin{itemize}
        \setlength{\itemsep}{0.7em}
        \item \textbf{Topology} - Dual Symmetrical Monopole with $\pm400kV$
        \item \textbf{DC Line} - DC cable and OHL can be utilized with the given current rating
        \item \textbf{Stages} - Two $1500MW$ stages
        \item \textbf{Flexibility} - Design can be used for configurations with $1000MW$ rating
      \end{itemize}
    \end{column}
  \end{columns}
\end{frame}


\begin{frame}{HVDC Configuration - BIP}
  \begin{columns}[c,onlytextwidth]
    \begin{column}{0.46\textwidth}
      \centering
      \topoimg[width=\linewidth]{chapters/system-configuration/bip}
    \end{column}
    \begin{column}{0.50\textwidth}
      \vspace{0.5em}
      \begin{itemize}
        \setlength{\itemsep}{0.7em}
        \item \textbf{Topology} - Bipole with $\pm525kV$
        \item \textbf{DC Line} - OHL can be utilized with the given current rating but parallel cables are required
        \item \textbf{Stages} - Two $1500MW$ stages
        \item \textbf{Flexibility} - Limited flexibility as asymmetrical configuration is not economical for most requirements
      \end{itemize}
    \end{column}
  \end{columns}
\end{frame}


\begin{frame}{HVDC Configuration - Rigid BIP}
  \begin{columns}[c,onlytextwidth]
    \begin{column}{0.46\textwidth}
      \centering
      %\topoimg[width=\linewidth]{chapters/system-configuration/rbip}
      \topoimg[width=\linewidth,height=0.85\textheight,keepaspectratio]{chapters/system-configuration/rbip}
    \end{column}
    \begin{column}{0.50\textwidth}
      \vspace{0.5em}
      \begin{itemize}
        \setlength{\itemsep}{0.7em}
        \item \textbf{Topology} - Dual Rigid Bipole with $\pm400kV$ and \textbf{3ph} transformer units
        \item \textbf{DC Line} - DC cable and OHL can be utilized with the given current rating
        \item \textbf{Stages} - $750MW$ stages for converter and $1500MW$ stages for DC line
        \item \textbf{Flexibility} - Limited flexibility as asymmetrical configuration is not economical for most requirements
      \end{itemize}
    \end{column}
  \end{columns}
\end{frame}


% \begin{frame}{Working Packages}
\begin{frame}{Working Packages Ongoing}
  \begin{itemize}
  \setlength{\itemsep}{0.5em}
  \item \textbf{DC line}
    \begin{itemize}
      \item Local company is working on DC route definition and DC line design

    \end{itemize}
  
    \item \textbf{Network Control \& Protection}
    \begin{itemize}
      \item Operation concept definition
      \item Offline studies for concept verification
      \item Development of models for performance design and verification
      \item Definition of performance verification tests and evaluation criteria
      \item Control hardware definition
    \end{itemize}
  
    % \item \textbf{HVDC concept evaluation}
    % \begin{itemize}
    %   \item Technical evaluation of HVDC concepts
    %   \item Commercial evaluation of HVDC concepts
    %   \item OEMs should be approached with only one concept
    % \end{itemize}
\end{itemize}
\end{frame}


% \begin{frame}{Working Packages}
%   \begin{itemize}
%   \setlength{\itemsep}{0.5em}
%   \item \textbf{Project execution}
%     \begin{itemize}
%       \item Converter specification to receive a binding offer (FEED phase avoided to reduce the execution time)
%       \item Supplier evaluation and selection (it is expected that suppliers from Asia might be more suitable for our approach)
%       \item Testlab buildout for real time simulations
%       \item Project execution support (Owners Engineer)
%     \end{itemize}
% \end{itemize}
% \end{frame}


\begin{frame}{Working Packages\\ HVDC configuration by build size and phasing}
  \begin{itemize}
  \setlength{\itemsep}{0.5em}
  \item Recommended configuration for $300MW$, $500MW$, $1GW$, $2GW$, $3GW$
    \begin{itemize}
      \item Monopole vs bipole by capacity
      \item Voltage class by size and distance: $\pm320kV$, $\pm400kV$, $\pm525kV$

    \end{itemize}
  
    \item Modular expansion path from first stage to GW scale
    \begin{itemize}
      \item Oversizing on day 1 vs future expansion
      \item Oversizing for redundancy
      \item Expansion from $300-500MW$ initial build to $1-2GW$
      \item Design approach for staged customer commitments
      \item Cost and schedule penalties from deferred expansion
      \item Avoiding stranded capex
    \end{itemize}  
\end{itemize}
\end{frame}


